\section{Consequences and usable specializations}
\label{sec:applications}

\subsection{How much the multivariable bound improves}
The earlier fixed-field recurrence was
\[
 M_1=h_0+\floor{eB},\qquad
 M_{j+1}=h_0+\floor{e\,\varphi(p^{M_j})B}.
\]
The repeated cyclotomic extensions are the source of its growth. Our substitutions use integer polynomials throughout and do not enlarge $K$. A single scalar specialization is performed at the end.

For $p=2$, $K=\Q_2$, $\rho=1$, $B=3$, and $h_0=1$, the closed bound in Theorem~\ref{thm:no-tower} simplifies to
\[
 H=2\cdot3^{d-1}+2.
\]
\begin{center}
\begin{tabular}{@{}rrr@{}}
\toprule
Variables $d$ & Previous recurrence & New sufficient $H$\\
\midrule
1 & $4$ & $4$\\
2 & $25$ & $8$\\
3 & $50{,}331{,}649$ & $20$\\
4 & $1+3\cdot2^{50{,}331{,}648}$ & $56$\\
\bottomrule
\end{tabular}
\end{center}
These compare sufficient bounds, not the true optimal endpoint. The scalar cost is sharp; the multivariable bound need not be. Computing the actual finite-slope radius recurrence and using the exact scalar tail can improve the closed bound further.

\subsection{A small example showing why one slope is insufficient}
Over $\Z_2[X,Y]$ take
\[
 F_n(X,Y)=Y-\big((1+X)^n-1\big),\qquad n=1,2,3.
\]
At $\rho=1$ all three Gauss valuations are $1$, and $(X,Y)=(0,2)$ is a common witness with all three valuations equal to $1$. The packet budget is $3$.

The substitutions $Y=(1+X)^n-1$ for $n=1,2,3$ each make one packet member identically zero. Taking a torsion parameter of larger order on any one of those curves cannot fix that failure. The finite-slope theorem considers four slopes and selects the fourth. If $\zeta$ is primitive of order $8$, the explicit tuple
\[
 X=\zeta-1,\qquad Y=\zeta^4-1=-2
\]
already works. Its three evaluated valuations are, respectively,
\[
 v(\zeta^4-\zeta)=\tfrac14,\qquad
 v(\zeta^4-\zeta^2)=\tfrac12,\qquad
 v(\zeta^4-\zeta^3)=\tfrac14.
\]
All are at most $1$, and both coordinates are nontrivial $2$-power torsion translates. The universal theorem allows the larger endpoint $H=8$; this example needs only $H=3$. It illustrates how the theorem preserves separate bounds even when several natural candidate curves fail completely.

\subsection{Integral formal power series}
The extension to $\OO_K[[T_1,\ldots,T_d]]$ is especially relevant to interpolation algebras. Integral coefficients ensure convergence at every tuple of positive coordinate valuations. Since the endpoint $H$ depends only on $p,e,d,\rho,B,h_0$, choose a total-degree cutoff $D$ with
\[
 (D+1)r_H>B.
\]
At any candidate tuple, every omitted monomial has valuation $>B$, uniformly over its degree. The same holds at the original radius $\rho$, because $\rho\ge r_{h_0}\ge r_H$. Truncation therefore preserves all $A_i\le B$, and after specialization the tail cannot cancel the selected value of valuation at most $A_i$. This justifies the series statement without a degree-dependent circularity.

The finite grid argument also permits an arbitrary family when every finite sum of its nonnegative Gauss budgets is at most $B$. This does not give a uniform bound for an arbitrary unbudgeted infinite family.

\subsection{Matrices and simultaneous denominator control}
\begin{corollary}[Preservation of selected minors]
\label{cor:minors}
Let $\mathcal M_i(\mathbf T)$ be finitely many matrices over
$\OO_K[[T_1,\ldots,T_d]]$. For each $i$ select a nonzero square minor $f_i$, and suppose its Gauss valuation $A_i$ is finite with $\sum_iA_i\le B$. Then the endpoint of Theorem~\ref{thm:no-tower} yields a common torsion tuple at which every selected minor is nonzero and
\[
 v(f_i(\boldsymbol\zeta-\mathbf1))\le A_i.
\]
In particular, every corresponding matrix has rank at least the size of its selected minor. If a chosen matrix is square, its inverse entries have valuation at least $-A_i$ at this tuple.
\end{corollary}
\begin{proof}
Apply the power-series theorem to the chosen determinants. Nonzero follows from a finite valuation bound. In the square case, the adjugate has integral entries at any positive-valuation tuple. Dividing by the determinant gives the asserted lower bound on the inverse entries.
\end{proof}
The choice of a nonzero minor is an input, not a conclusion that can be obtained from a generic determinant bound. Witness evaluations with finite valuations may be used to provide the needed budgets. This corollary sharpens the order cutoff in the earlier matrix/Fitting applications; it does not claim those applications themselves are new.

\subsection{Quantifier discipline in an arithmetic application}
The new interface is
\[
 \text{fix }K,p,d,\rho,B
 \quad\Longrightarrow\quad
 \text{choose }H
 \quad\Longrightarrow\quad
 \text{handle every admissible packet over }\OO_K.
\]
The packet can have arbitrary degrees and arbitrarily many members. Its coefficients must nevertheless be defined over the fixed field, or over fields with a known uniform bound on the wild ramification exponent.

By contrast, the upstream uniform selector can handle arbitrary algebraic integral coefficients after its order bound is chosen, at positive tolerance. A finite packet always lies in some finite field, but selecting that field \emph{after} inspecting the packet does not supply a uniform bound for the preceding quantifier order. Consequently our theorem is a valid replacement in fixed-field interpolation settings, not a blanket removal of the tolerance from every upstream Selmer argument.

Even a perfectly controlled interpolation polynomial does not by itself yield nonvanishing of an arithmetic special value. One must first establish that the polynomial represents the desired object, that the integral normalization is correct, and that its baseline is finite. Those are the substantive arithmetic inputs.

\subsection{What remains open locally}
The results leave three clear refinements:
\begin{enumerate}
\item Determine the positive-tolerance obstruction cost for general noncyclotomic coefficient fields. The zero-tolerance formula does not answer that more refined question.
\item Determine the optimal dependence on dimension in the multivariable bound. The factor $3^{d-1}$ comes from a convenient interpolation estimate, not an optimality argument.
\item Determine sharp bounds under restrictions on packet cardinality or on actual arithmetic interpolation families. An arbitrary polynomial obstruction need not occur as a Selmer determinant.
\end{enumerate}
These are now more focused questions than whether the earlier field-extension tower is necessary: the latter is answered negatively by Theorem~\ref{thm:no-tower}.

